% !TEX root = ../notes.tex

\documentclass[../notes.tex]{subfiles}

\begin{document}

\section{April 1}
Here we go.

\subsection{The Local-Global Principle}
We are interested in finding rational points on varieties. Let's start with some notation for the present subsection. Fix a global field $k$, like $\QQ$ or $\FF_q(t)$, and for each place $v$, we let $k_v$ be the corresponding complete local field. If $v$ is finite, we let $\OO_v$ be its ring of integers; we may as well define $\OO_v=k_v$ when $v$ is archimedean. As usual, we define the ring of adeles $\AA_k$ to be the restricted product
\[\AA_k\coloneqq\prod_v(F_v,\OO_v),\]
which we recall is locally compact.

Merely in terms of finding rational points, it turns out that one can reduce to the case where the variety is smooth, projective and geometrically irreducible.
\begin{definition}[nice]
	Fix a variety $X$ over a field $k$. Then $X$ is \textit{nice} if and only if $X$ is smooth, projective, and geometrically irreducible.
\end{definition}
\begin{remark}
	Given smoothness, one can weaken ``geometrically irreducible'' to ``connected.''
\end{remark}
\begin{remark}
	It turns out that nice varieties admit $k_v$-rational points for all but finitely many $v$, and this finite set can be made effective.
\end{remark}
In general, there is no known algorithm to determine if a variety $X$ over a global field $k$ admits a rational point. However, there are some known cases.
\begin{theorem}[Hasse--Minkowski] \label{thm:hasse-minkowski}
	Fix a global field $k$, and let $X$ be a nice quadric hypersurface in $\PP^n_k$. Then $X(k)$ is nonempty if and only if $X(k_v)$ is nonempty for all places $v$.
\end{theorem}
\begin{remark}
	This was proven in the 1890s over $\QQ$ by Minkowski and for general number fields by Hasse in 1924.
\end{remark}
The moral is that one can test if $X(k_v)$ is nonempty for each place $v$ fairly explicitly: there is no content over $\CC$, one uses some inequalities over $\RR$, and one can use Hensel's lemma for finite places.
% \begin{remark}
% 	It turns out that even in $\PP^2$, being able to find the rational point is at least as hard as integer factorization.
% \end{remark}

The conclusion of \Cref{thm:hasse-minkowski} is nice enough to be worth a name.
\begin{definition}[Hasse principle]
	A variety $X$ over a global field $k$ satisfies the \textit{Hasse principle} or the \textit{local-global principle} if and only if $X(k)$ is nonempty is equivalent to $X(k_v)$ being nonempty for every $v$.
\end{definition}
\begin{remark}
	For projective $X$, one has that $X(\OO_v)=X(k_v)$ by the valuative criterion for properness, so
	\[X(\AA_k)=\prod_vX(k_v).\]
	Indeed, one just has to choose a model for $X$ which is projective away from finitely many places. The moral is that the Hasse principle can be stated as having $X(k)\ne\emp$ equivalent to $X(\AA_k)\ne\emp$.
\end{remark}
The Hasse principle can fail.
\begin{example}
	Hasse showed in 1931 that the variety of norms from $\QQ(\sqrt{-3},\sqrt{13})$ fails the Hasse principle.
\end{example}
\begin{example}
	Lind and Reichardt in the 1940s produced some examples of genus $1$ curves where the Hasse principle fails. Some examples of such curves are as $3x^3+4y^2+5z^3=0$ (in $\PP^2_\QQ$) and a smooth projective model of $2y^2=1-17x^4$.
\end{example}
\begin{remark}
	For elliptic curves $E$ over a global field $k$, one has that $E$ has a nonzero point over every local field (by some argument with the Implicit function theorem), but there is no reason for $E$ to have a nonzero rational point.
\end{remark}
\begin{example}
	Swinnerton-Dyer constructed some cubic surfaces where the Hasse principle fails. Later, Cassels and Guy gave a ``diagonal'' example $5x^3+9y^3+10z^3+12w^3=0$ in $\PP^3_\QQ$.
\end{example}
In 1970, Manin would explain what is going on in all these examples, which is the point of the present course.

\subsection{The Brauer Group, from Algebra}
Let's start with some algebra. Throughout, when $k$ is a field, we let $\overline k$ denote an algebraic closure, and we let $k^{\mathrm{sep}}$ be the separable closure of $k$ in $\overline k$. For convenience, we also set $G_k\coloneqq\op{Gal}(k^{\mathrm{sep}}/k)$.
\begin{definition}
	An \textit{algebra} over a field $k$ is an associative (unital) ring $A$ with a ring homomorphism $k\to A$.
\end{definition}
One has the following structure theorem for so-called ``central simple'' algebras.
\begin{proposition} \label{prop:csa-grab-bag}
	Fix a finite-dimensional $k$-algebra $A$. Then the following are equivalent.
	\begin{listroman}
		\item One has that $A\otimes_kk^{\mathrm{sep}}\cong M_n(k^{\mathrm{sep}})$ for some $n\ge1$.
		\item One has that $A\otimes_k\overline k\cong M_n(\overline k)$ for some $n\ge1$.
		\item The algebra $A$ is central (i.e., $Z(A)=k$) and simple (i.e., $A$ admits exactly two ideals).
		\item One has that $A\cong M_r(D)$ for some $r\ge1$ and some finite-dimensional central division $k$-algebra $D$.
	\end{listroman}
\end{proposition}
\begin{proof}
	Omitted. This is a standard result in non-commutative algebra.
\end{proof}
\begin{remark}
	In (iv), it turns out that $r$ and $D$ are uniquely determined by $A$.
\end{remark}
\begin{example} \label{ex:real-csa}
	Over $\RR$, the only central division algebras are $\RR$ and $\HH$, so the central simple algebras look like $M_r(\RR)$ and $M_r(\HH)$.
\end{example}
\begin{example}[quaternion algebra] \label{ex:quaternion-csa}
	For any field $k$ of characteristic not $2$, one can define a quaternion algebra as follows: choose $a,b\in k^\times$, and then
	\[H_{a,b}\coloneqq\frac{k\langle i,j\rangle}{\left(i^2-a,j^2-b,ij+ji\right)}.\]
	By reducing words, one sees that $H_{a,b}$ is four-dimensional (spanned by $\{1,i,j,ij\}$), and one can check that this is central simple. Note that this does not have to be a division algebra: for example, if $a$ is a square, then $H_{a,b}$ admits zero-divisors because $(i-\sqrt a)(i+\sqrt a)=0$.
\end{example}
\begin{remark}
	One can check that all four-dimensional central simple algebras $A$ are isomorphic to some quaternion algebra. As a sketch, one finds some $i\in A\setminus k$, which one can check must satisfy some quadratic equation. By completing the square, one may assume that $i^2=a$ for some $a\in k^\times$. Then one can diagonalize the action of $i$ in $A$ (by conjugation), and choose $j$ to be in the $(-1)$-eigenspace.
\end{remark}
Our structure theory allows us to partition central simple algebras.
\begin{definition}[Brauer group]
	Fix a field $k$. Two central simple algebras $A$ and $B$ are \textit{equivalent} if and only if there is a division algebra $D$ and integers $r$ and $s$ such that $A\cong M_r(D)$ and $B\cong M_s(D)$. The \textit{Brauer group} $\op{Br}(k)$ is the collection of central simple $k$-algebras, considered up to equivalence.
\end{definition}
\begin{remark} \label{rem:better-equiv}
	We will not bother to check that equivalence is an actual equivalence relation. One can remove the division algebra $D$ by instead asking for $M_s(A)\cong M_r(B)$. Indeed, the point is that
	\[M_r(M_s(D))=M_{rs}(D)=M_s(M_r(D)).\]
\end{remark}
\begin{remark} \label{rem:br-is-group}
	The collection $\op{Br}(k)$ is a group under $\otimes$ (which turns out to send central simple algebras to central simple algebras). The unit is $k$, and the inverse of $A$ is $A\opp$; for example, one needs to check that $A\otimes A\opp$ is isomorphic to a matrix algebra over $k$ (it's $\op{End}_k(A)$).
\end{remark}
\begin{remark}
	A field extension $k\into\ell$ induces a group homomorphism $\op{Br}(k)\to\op{Br}(\ell)$ given by $-\otimes_k\ell$. Notably, $-\otimes_k\ell$ sends central simple algebras to central simple algebras by \Cref{prop:csa-grab-bag}(ii), and it preserves equivalence by the definition given in \Cref{rem:better-equiv}.
\end{remark}

\subsection{The Brauer Group, from Cohomology}
There is also a cohomological construction of the Brauer group.
\begin{proposition} \label{prop:brauer-cohomology}
	Fix a field $k$. Then $\op{Br}(k)\cong\mathrm H^2(k;k^{\mathrm{sep}\times})$.
\end{proposition}
\begin{proof}
	For each $n$, \Cref{prop:csa-grab-bag}(i) says that central simple algebras of dimension $n^2$ are Galois twists of $M_n(k)$, so they are classified by the pointed set $\mathrm H^1(k;\op{Aut}M_n(k^{\mathrm{sep}}))$. By the Skolem--Noether theorem, all (algebra) automorphisms of $M_n(k^{\mathrm{sep}})$ are inner, so this automorphism group is $\op{PGL}_n(k^{\mathrm{sep}})$. To compute this cohomology group, we note that there is a short exact sequence
	\[1\to\mathbb G_m\to\mathrm{GL}_n\to\mathrm{PGL}_n\to1\]
	of group schemes, so taking the long exact sequence in cohomology, we receive an exact sequence
	\[\mathrm H^1(k;\mathrm{GL}_n)\to\mathrm H^1(k;\mathrm{PGL}_n)\to\mathrm H^2(k;\mathbb G_m)\]
	of pointed sets. We have $\mathrm H^1(k;\mathrm{GL}_n)=\{1\}$ by (a version of) Hilbert's theorem 90. (This is basically the fact that every $n$-dimensional vector space over a Galois extension of $k$ with a semilinear action of the Galois group comes from an $n$-dimensional vector space over $k$.) It turns out that this is enough to produce an injection $\mathrm H^1(k;\mathrm{PGL}_n)\to\mathrm H^2(k;\mathbb G_m)$.\footnote{Technically, because this is nonabelian cohomology, injectivity requires $\mathrm H^1(k; (A\otimes_k k^{\mathrm{sep}})^\times)$ to be trivial for every central simple algebra $A$ over $k$ of dimension $n^2$, but this too is true.}

	Now, taking the colimit over $n$, one shows that the collection of all central simple algebras to $\mathrm H^2(k;\mathbb G_m)$. It remains to check that this map descends to equivalence classes and then defines an isomorphism on $\op{Br}(k)$. We will not run any of these checks.
\end{proof}
\begin{remark}
	When $\op{char}k$ does not divide $n$, one sees that the map $\op{SL}_n(k^{\mathrm{sep}})\to\op{PGL}_n(k^{\mathrm{sep}})$ is surjective with kernel $\mu_n$. Thus, each element in the image of the natural map $\mathrm H^1(k;\mathrm{PGL}_n)\to\mathrm H^2(k;\mathbb G_m)$ is also the image of some element of $\mathrm H^2(k;\mu_n)$, which is $n$-torsion. The same argument works in characteristic dividing $n$ if we use flat cohomology instead of \'etale (Galois) cohomology; namely, this allows us to pass from $k^{\mathrm{sep}}$ to $\overline k$ everywhere. It follows that $\op{Br}(k)$ is a torsion group.
\end{remark}
Now that we know that $\op{Br}(k)$ is torsion, we can hope to compute its $n$-torsion. This is a variant of Kummer theory.
\begin{proposition} \label{prop:kummer}
	Fix a field $k$ of characteristic not dividing some positive integer $n$. Then $\mathrm H^1(k;\mu_n)=k^\times/k^{\times n}$ and $\mathrm H^2(k;\mu_n)=\op{Br}(k)[n]$.
\end{proposition}
\begin{proof}
	We use the short exact sequence
	\[1\to\mu_n\to\mathbb G_m\stackrel n\to\mathbb G_m\to1\]
	of \'etale group schemes on $k$. Taking Galois cohomology, we receive a long exact sequence
	\[k^\times\stackrel n\to k^\times\to\mathrm H^1(k;\mu_n)\to\mathrm H^1(k;\mathbb G_m).\]
	The last group vanishes by Hilbert's theorem 90, so $\mathrm H^1(k;\mu_n)\cong k^\times/k^{\times n}$ follows. Similarly, we have a long exact sequence
	\[\mathrm H^1(k;\mathbb G_m)\to\mathrm H^2(k;\mu_n)\to\mathrm H^2(k;\mathbb G_m)\stackrel n\to\mathrm H^2(k;\mathbb G_m),\]
	and the left group vanishes by Hilbert's theorem 90, so $\mathrm H^2(k;\mu_n)\cong\mathrm{Br}(k)[n]$ follows by \Cref{prop:brauer-cohomology}.
\end{proof}
\begin{remark}
	If we use flat cohomology, then the same proof works even if $\op{char}k\mid n$, so the conclusion is recovered.
\end{remark}
Having a cohomological interpretation of the Brauer group allows us to do Galois cohomology things to it. Here is an example calculation.
\begin{example} \label{ex:quaternion-by-cup}
	Suppose that $\op{char}k$ is not $2$. There is a canonical cup product
	\[\mathrm H^1(k;\mu_2)\times\mathrm H^1(k;\ZZ/2\ZZ)\to\mathrm H^2(k;\mu_2).\]
	Of course, there is a canonical isomorphism $\mathrm H^1(k;\ZZ/2\ZZ)\cong\mathrm H^1(k;\mu_2)$, so we see that this induces a cup product
	\[k^\times/k^{\times2}\times k^\times/k^{\times2}\to\op{Br}(k)[2]\]
	by \Cref{prop:kummer}. One can check that this takes a pair $(a,b)$ to the quaternion algebra $H_{a,b}$, which we can do with some explicit cocycle calculations to $\mathrm H^2(k;\mu_2)$.
\end{example}
\begin{remark}
	For fixed $a\in k^\times$, it turns out that the natural map $k^\times\to\op{Br}(k)$ given by $b\mapsto[H_{a,b}]$ has kernel given by $\op N_{k(\sqrt a)/k}k(\sqrt a)^\times$. Indeed, the idea is to use Tate Galois cohomology: there is an explicit isomorphism
	\[\mathrm H^0_T(\op{Gal}(k(\sqrt a)/k);\mathbb G_m)\to\mathrm H^2(\op{Gal}(k(\sqrt a)/k);\mathbb G_m)\]
	given by some cup product, computed in \Cref{ex:quaternion-by-cup}. The source reduces to $k^\times/\op N_{k(\sqrt a)/k}k(\sqrt a)^\times$ by expanding out the Tate cohomology group, and the target can be seen to be the kernel of the natural map $\op{Br}(k)\to\op{Br}(k(\sqrt a))$ by the Inflation--restriction exact sequence.
\end{remark}

\subsection{Examples of Brauer Groups}
% It is as good a time as any to give some examples of Brauer groups.
\begin{example}
	If $k$ is separably closed, then \Cref{prop:csa-grab-bag} shows that $\op{Br}(k)$ is trivial. One can also see this directly from \Cref{prop:brauer-cohomology}.
\end{example}
\begin{example}
	By Wedderburn's theorem, all finite division algebras are fields, so there is only one finite central division algebra over $\FF_q$ (namely, $\FF_q$). Thus, $\op{Br}\FF_q=0$.
\end{example}
\begin{example}
	By Tsen's theorem, one has that $\op{Br}\CC(t)=0$. In fact, the same is true for any field of transcendence degree $1$ over an algebraically closed field.
\end{example}
\begin{example}
	A direct calculation in Galois cohomology shows that $\op{Br}\RR=\frac12\ZZ/\ZZ$, where $\RR$ goes to $0$, and $\HH$ goes to $\frac12$. This is equivalent to \Cref{ex:real-csa}.
\end{example}
\begin{example}
	For any nonarchimedean local field $k_v$, it turns out that $\op{Br}k_v\cong\QQ/\ZZ$. There is a way to write down these elements explicitly. This assertion is roughly as hard as local class field theory.
\end{example}
\begin{example} \label{ex:cft-exact-seq}
	For any global field $k$, functoriality gives a map $\op{Br}(k)\to\prod_v\op{Br}k_v$. It turns out that the image lands in $\bigoplus_v\op{Br}k_v$, which one sees because a central simple algebra over $k$ uses finitely many denominators, so it can be trivialized at all but finitely many places $v$. By global class field theory, it turns out that this map is injective, and it fits into a short exact sequence
	\[0\to\op{Br}k\to\bigoplus_v\op{Br}k_v\to\QQ/\ZZ\to0,\]
	where the last map is given by summing. Similarly, this is roughly as hard as global class field theory.
\end{example}
\begin{remark}
	If $p$ and $q$ are distinct odd primes, then the fact that $H_{p,q}$ goes to zero under the composition $\op{Br}\QQ\to\bigoplus_v\op{Br}\QQ_v\to\QQ/\ZZ$ is equivalent to quadratic reciprocity.
\end{remark}

\end{document}